\begin{pastexam}[label={px:nchu-113-dsa-18}]{中興資工所甲組 113 學年度 選擇題第 18 題}
（單選；答錯不倒扣）

A robot is located at the top-left corner of an m$\times$n grid. The robot can only move either down or
right at any point in time. The robot is trying to reach the bottom-right corner of the grid. How many
possible unique paths are there for a 3$\times$7 grid, using a dynamic programming approach? (5\%)
\begin{enumerate}[label=(\Alph*)]
\item 21
\item 28
\item 35
\item 42
\end{enumerate}
\end{pastexam}

\begin{pastsoln}
答案是 (B)。

把格子看成 $3$ 列、$7$ 行，機器人從左上角那一格走到右下角那一格。 令 $P[i][j]$ 是走到第 $i$ 列第 $j$ 行
（從 $1$ 數起）的路徑數。 最後一步不是從上面來就是從左邊來，兩類不重疊，所以
\[
P[i][j] = P[i-1][j] + P[i][j-1],\qquad P[1][j] = P[i][1] = 1
\]
（第一列只能一路往右、第一行只能一路往下，各只有一種走法）。 由上而下、由左而右填表（資結書第~18.1.3~節
的由下往上）：
\begin{center}
\begin{tabular}{c|ccccccc}
 & 1 & 2 & 3 & 4 & 5 & 6 & 7 \\
\hline
1 & 1 & 1 & 1 & 1 & 1 & 1 & 1 \\
2 & 1 & 2 & 3 & 4 & 5 & 6 & 7 \\
3 & 1 & 3 & 6 & 10 & 15 & 21 & \textbf{28} \\
\end{tabular}
\end{center}
$P[3][7] = 28$，選 (B)。 檢查：每條路徑恰好往下 $2$ 步、往右 $6$ 步，路徑數是在 $8$ 步裡選哪 $2$ 步往下，
$\binom{8}{2} = 28$。 若把題目讀成 $3 \times 7$ 個格點之間的線段（往下 $3$、往右 $7$），會得到
$\binom{10}{3} = 120$，不在選項裡，可見題目的 grid 指的是格子。
\end{pastsoln}
